% NAL 2338, batch 3 — Gallica views 41–60.
% Pass: Opus 5.5 (claude-opus-5-5), 2026-09-23 — first pass, unchecked against the views by a human.
% Read from the local mirror only (archives/nal-2338/f0041.jpg … f0060.jpg):
% 1000-px thumbnails first, then four full-width bands per written view at
% 2400 px (six bands for the denser hand of v. 57–60), then tight crops for
% doubtful words and for the faded patches of v. 54 and v. 57. Each view was
% drafted to the scratchpad (nal-2338/drafts-b3/) as it was read, and this
% file assembled from the drafts. No edition of Torricelli, Roberval or
% Mersenne and no other transcription was opened.
%
% Dating. The catalogue gives none: \dating{} is left empty, and no date is
% inferred for the volume. The leaves carry their own dates, transcribed where
% they stand: « Dat. Florentiae die 7 Julij ann. 1646 » (v. 50) and « flor. Die
% [8?] Julij an. 1646 » (v. 55).
%
% What the twenty views hold. Colour and greyscale scans, portrait, about
% 4316 × 5750, one page per view with the edge of the next leaf or the binding
% at one side. Three Latin letters, all copies in fair hands (no leaf in this
% batch is in the hand of the man it names as writer, as far as the leaves
% show; they are copies, not autographs, by their layout and regularity —
% an observation, not an attribution).
%
%   v. 41–50   Letter 1, continued from before v. 41: v. 41 opens mid-sentence
%              inside a parenthesis (« intentum meum) ipsius solidum
%              demonstrare »). Its heading is on the recto before (batch 2),
%              seen reversed through v. 41 (« … Roberuallio … »). The
%              addressee is Roberval by the text's own words (« linea
%              Roberualliana … quam ipse excogitasti Vir Clar.me », v. 47;
%              « linea tua », v. 48). Priority over the centre of gravity of
%              the cycloid, the tangents « ex doctrina motus », Mersenne's
%              letter quoted (v. 42), the writer's « libellus de Motu » and
%              Galileo's book on projectiles (v. 43–44), Guldin, « Caua[?] »
%              (v. 44), Desargues (« de sargues », v. 44); then, from v. 45,
%              Lemma I (tangent to any parabola, exponents 5 and 3 as
%              example), Lemma II (parallelograms to inscribed parabolas of
%              one species), two quadratures of all parabolas (v. 46–48), the
%              « quinque methodi » for tangents, a spiral « ex infinitis
%              spiris conflata » equal to a straight line (v. 49), hyperbola
%              theorems sent to Fermat (v. 49). Ends v. 50: « Vale. Dat.
%              Florentiae die 7 Julij ann. 1646 », then a postscript saying
%              the letter was held until the beginning of August. Unsigned;
%              the writer is not named in views 41–60.
%   v. 51, 52  The same blank verso (of leaf 25) photographed twice, with
%              show-through. Not transcribed.
%   v. 53–55   Letter 2, headed « Ill.mo et Doctissimo Viro P. de Carcauy
%              Euang. Torricellius Sal. », same copyist as letter 1: Fermat's
%              numerical problems declined, the problem of the infinite
%              parabolas (Roberval, Fermat, Cavalieri 1639, Niceron),
%              methods for tangents and quadratures, a rule for centres of
%              gravity stated with « omnes &c. ad omnes &c. », the
%              « Aristarchi libellus », a request that the infinite
%              hyperbolas and the spirals be made known to Fermat. Signed in
%              the text « … ama Euangelistam Torricellium. D. flor. Die [8?]
%              Julij an. 1646. », address « Ill.mo Viro P. de Carcauy » at the
%              foot of v. 55. Begins and ends within the batch. The ink of v. 54
%              has faded on four lines (≈ 8 \ill{}).
%   v. 56      Blank verso (of leaf 27). Not transcribed.
%   v. 57–60   Letter 3, headed « Epistola Aegidii Personerii de Roberval ad
%              [struck: R. P.] Euangelistam Torricellium », above two struck
%              French lines « Lettre de Mr De Roberval escrite au Sre
%              Toricelli »; a second hand, smaller and closer, with accents.
%              « Vir clarissime »: an answer to charges of plagiarism; the
%              trochoid first heard of from Mersenne in Paris in 1628, left
%              six years, taken up in 1634; the « infiniti doctrina » against
%              Cavalieri's indivisibles; tangents by composition of motions
%              « circa annum 1636 », lectures collected by du Verdus;
%              correspondence with Fermat through Carcavy from 1635. Runs on
%              past v. 60 in mid-sentence (« imò vniuersa-/lius »). No date in
%              these views. A French note at the foot of v. 57, in another
%              hand, tells someone to set « de caracteres Italiques » what is
%              in large characters here — a printer's or copyist's
%              instruction; set as \marginal{} with a note.
%
% First and last views. v. 41 continues a letter begun before it (the
% heading of letter 1 is on the recto before, in batch 2). v. 60 runs on into
% v. 61 in mid-word; letter 3 is not finished in this batch.
%
% Hands. (A) v. 41–55: a regular copyist's italic with large ornamental
% initials and long line-fillers at line ends (fillers, not deletions);
% « q; » for -que, « gr. » for gravitatis, a barred p for per (set « p. »),
% « et. c. » for etc.; its r and v are close (« vatio » is « ratio », v. 53).
% (B) v. 57–60: smaller and denser, accents on vowels (à, â, ò, ù, è),
% question marks as a hook over a point, dates with a hooked 1 and a b-shaped
% 6 (« ıb36 » = 1636). (C) the French heading and foot-note of v. 57, paler.
%
% Spelling. The copies' own: u inside words and v at their head (« Inuersio »,
% « vnicuiq; »), ij, « quodamodo », « dilibantur », « laberintos », « Caua »,
% « de sargues », « De Ferma ». A long s is transcribed as s. Accents and the
% marks of abbreviation over a vowel (lineā, quandā, tamē, sūpto) are kept as
% on the page. Latin as written; not normalised.
%
% Notation. Points of figures in capitals in the running text (ABC, B.I, the
% dots are the copy's); no figure is drawn anywhere in these views, though the
% lemmas of v. 45–48 suppose them and leave half the page empty for them. No
% sign of operation: « cuboquadratum », « ductus … in … » in words.
%
% Numbers on the leaves.
% - One ink series at the top right of the rectos: 21 (v. 42), 22 (v. 44),
%   23 (v. 46; its top also shows above v. 44 on the next leaf), 24 (v. 48),
%   25 (v. 50), 26 (v. 53), 27 (v. 55), 28 (v. 57), 29 (v. 59); « 20 » and
%   « 21 », « 22 » seen reversed through the versos v. 41, 43, 45. It runs on
%   across hands A and B and counts the leaves whose versos are blank (25,
%   27), so it behaves as the library's foliation; but it is ink, so, as for
%   NAF 5176, NAF 9544 and Français 12357, no \folio{} is written until a
%   human has looked. No certificate leaf is in this batch to check it
%   against. No verso carries a number.
% - Dates and years in the text are content: 1628, 1634, 1635, 1636, 1639,
%   1642, 1646; « 27um » (Roberval's age, v. 58); « 7 ad 5 » (v. 42);
%   exponents 5, 3, « 3 ad 2 » (v. 45–46); « 19 Quinti » (v. 47).
% - A small « z » or « 2 » in ink at the right edge of v. 44, meaning not
%   established.
% - Stamps: the round « BIBLIOTHÈQUE NATIONALE MSS » with « R.F. » (v. 53, 57).
%
% The catalogue's dating is copied verbatim (empty).

% Coordinator's reconciliation (2026-09-23), after all six batches:
%   One ink series of leaf numbers runs top right of the rectos across every
%   hand and piece: 1–10 (v. 2–20), 11–20 (v. 22–40), 21–29 (v. 42–59),
%   30–38 with « 34 bis » for the figure slip (v. 61–79), 39–48 (v. 81–99),
%   49–50 (v. 101–103). It ends on « 50 » at the last written leaf, as the
%   certificate of v. 1 has it (« Volume de 50 Feuillets / 1er Octobre
%   1888 »): it is the library's count. Being ink, it gets no \folio{}, in any
%   batch. The bold series on the letter of v. 24–35 (1–11) is that letter's
%   own pagination (batch 2).
%   Seams: v. 60 ends « vniuersa- » and v. 61 opens « lius »; batch 3 had set
%   the catchword into v. 60, and now ends where the page does. v. 100 runs
%   into v. 101 (« altius ; » / « in aliis autem humilius »), the « De Vacuo
%   Narratio » (v. 97–103). The « non » sign is set « nl » in batch 1 (hand B,
%   an n with a barred ascender) and « n^o » in batch 2 (the fast cursive, an
%   n with a raised loop): two hands, two signs, left as each pass read them.

\documentclass[11pt,a4paper]{article}
% The preamble shared by every transcription of the mathematicians' manuscripts in Gallica.
%
% It rests on one idea: **uncertainty compounds, it does not resolve.** A
% transcription that smooths over a doubtful word has destroyed the only thing
% it was made for — a reader must be able to tell, at every point, what was on
% the page and what was supplied. The apparatus macros below are therefore the
% critical apparatus, not decoration.
%
% The same commands are recognised by `scripts/render.mjs`, which produces the
% site's reading view, and by `scripts/tei.mjs`. A macro added here must be
% added there too, or rendering fails loudly — which is the wanted behaviour.
%
% Comments here are English, like the rest of the repository. The documents
% this preamble sets are French, and that is deliberate: see the skills.

\ProvidesPackage{germain}[2026/09/15 Transcriptions of mathematicians' manuscripts in Gallica]

\usepackage{amsmath,amssymb,amsthm}
\usepackage{mathrsfs}
% Kept from the parent preamble so the renderer's diagram support stays
% available; these manuscripts draw no commutative diagrams, and nothing here uses it.
\usepackage{tikz-cd}
\usepackage[english,french]{babel}
\usepackage{fontspec}

% --- Greek ------------------------------------------------------------------
% The text face stays fontspec's default, Latin Modern, which has no Greek:
% XeTeX drops every glyph it lacks with a line in the log and nothing on the
% page, and the Greek words the manuscripts quote vanished from the PDFs. So
% the Greek and Greek Extended blocks get a XeTeX character class of their own,
% and entering or leaving a run of it switches the family to CMU Serif — the
% same Computer Modern design, with polytonic Greek — and back. Only the
% family changes: a Greek word in \textit or \textbf keeps its shape and series.
% The transitions to and from every other class carry the tokens that class
% has towards an ordinary letter, so French spacing before « ; » or inside
% « » still holds after a Greek word. No grouping, so no brace or \endgroup
% can come out unbalanced; maths is left alone, since XeTeX inserts nothing
% there. Kept identical in `exercices.sty`.
\makeatletter
\newfontfamily\germainGreekfont{cmun}[
  Extension=.otf, NFSSFamily=germaingreek,
  UprightFont=*rm, ItalicFont=*ti, SlantedFont=*sl,
  BoldFont=*bx, BoldItalicFont=*bi, BoldSlantedFont=*bl]
\newXeTeXintercharclass\germain@greek
\def\germain@greekrange#1#2{%
  \@tempcnta=#1\relax
  \loop
    \XeTeXcharclass\@tempcnta=\germain@greek
    \ifnum\@tempcnta<#2\advance\@tempcnta\@ne\repeat}
\germain@greekrange{"0370}{"03FF}
\germain@greekrange{"1F00}{"1FFF}
\def\germain@greekprev{\rmdefault}
\def\germain@greekin{%
  \edef\germain@greekprev{\f@family}\fontfamily{germaingreek}\selectfont}
\def\germain@greekout{\fontfamily{\germain@greekprev}\selectfont}
\def\germain@greekjoin{%
  \XeTeXinterchartokenstate=\@ne
  \@tempcnta=\z@
  \loop
    \ifnum\@tempcnta=\germain@greek\else
      \XeTeXinterchartoks\germain@greek\@tempcnta=\expandafter{%
        \expandafter\germain@greekout\the\XeTeXinterchartoks\z@\@tempcnta}%
      \XeTeXinterchartoks\@tempcnta\germain@greek=\expandafter{%
        \the\XeTeXinterchartoks\@tempcnta\z@\germain@greekin}%
    \fi
    \ifnum\@tempcnta<4095\advance\@tempcnta\@ne\repeat}
% babel-french sets its own classes, and saves and restores the interchar
% state, each time a language is selected: join again after it does.
\addto\extrasfrench{\germain@greekjoin}
\addto\extrasenglish{\germain@greekjoin}
\AtBeginDocument{\germain@greekjoin}
\makeatother

\usepackage{geometry}
\geometry{a4paper,margin=25mm,marginparwidth=28mm}
\usepackage{marginnote}
\usepackage{xcolor}
\usepackage[normalem]{ulem}
\setlength{\emergencystretch}{3em}
\usepackage{hyperref}
\hypersetup{colorlinks=true,linkcolor=black,urlcolor=black,pdfborder={0 0 0}}

% --- Batch metadata ---------------------------------------------------------
% Taken as they stand by the rendering script, which shows them at the head of
% the reading view. They fix what the file claims to cover. The macro names are
% the parent project's, so that the scripts stay shared; \folder here means the
% volume's slug — `fr-9115` — and \pages counts Gallica views.

\newcommand{\folder}[1]{\gdef\grothFolder{#1}}
\newcommand{\batch}[1]{\gdef\grothBatch{#1}}
\newcommand{\pages}[2]{\gdef\grothFirst{#1}\gdef\grothLast{#2}}
\newcommand{\foldertitle}[1]{\gdef\grothFoldertitle{#1}}

% The holder's own shelfmark, printed as they write it: « Français 9115 ».
\newcommand{\shelfmark}[1]{\gdef\grothShelfmark{#1}}

% The dating, copied verbatim from the holder's catalogue — never inferred
% here. « XIXe siècle » is the BnF's; a pass that dates a leaf from its content
% says so in a \note{}, as its own inference.
\newcommand{\dating}[1]{\gdef\grothDating{#1}}

% The status notice, printed in the title block and repeated as a diagonal
% watermark on every page of the PDF. The manuscripts are in the public
% domain; what this line declares is not a right but a quality — an
% unverified first pass by a machine. The skills write the line; a file
% without it gets no watermark, so leaving it out is visible in the output.
\newcommand{\watermark}[1]{\gdef\grothWatermark{#1}}

\gdef\grothFolder{?}\gdef\grothBatch{}\gdef\grothFirst{?}\gdef\grothLast{?}%
\gdef\grothFoldertitle{}\gdef\grothShelfmark{}\gdef\grothDating{}\gdef\grothWatermark{}

% --- The résumé -------------------------------------------------------------
\newenvironment{resume}
  {\small\setlength{\parskip}{0.45em}}
  {\par\medskip\hrule\medskip}

% --- Keywords ---------------------------------------------------------------
% In English, closing the modernised reading's résumé: the modern vocabulary
% under which to search for what the leaves construct. The manifest extracts
% them to tag the volume — this line is the single source of the tags.
\newcommand{\keywords}[1]{\par\smallskip\noindent
  {\footnotesize\textsc{Keywords} --- \textit{#1}}\par}

% --- The critical apparatus -------------------------------------------------

% The Gallica view number, in the margin. This is the anchor that lets a
% disagreement over a formula be settled by looking at *one* image rather than
% a volume of 750. The site uses it to turn the facsimile as the transcription
% is scrolled. A view is one scanned image: usually one side of a leaf.
\newcommand{\page}[1]{%
  \marginnote{\footnotesize\color{gray!70}\textsf{v.\,#1}}%
  \label{page:#1}\ignorespaces}

% The BnF's pencilled foliation, where the leaf carries it and a pass can read
% it: « 198r ». This is what the literature cites — « ff. 198r–208v » — and
% the one line that turns a citation into a view. Recorded, never inferred:
% a folio a pass cannot read is not written.
\newcommand{\folio}[1]{%
  \marginnote{\footnotesize\color{gray!50}\textsf{f.\,#1}}\ignorespaces}

% The modernised reading's coarser counterpart to \page{}: which views a
% section is drawn from.
\newcommand{\pagerange}[2]{%
  \marginnote{\footnotesize\color{gray!70}\textsf{v.\,#1--#2}}%
  \ignorespaces}

% Illegible: never guessed.
\newcommand{\ill}{\textcolor{red!60!black}{[\,\ldots\,]}}

% A doubtful reading: the word is given, and flagged as uncertain.
\newcommand{\uncertain}[1]{\uline{#1}}

% An editorial addition — a missing letter, an implied word.
\newcommand{\add}[1]{\textcolor{blue!50!black}{[#1]}}

% Struck out by the writer. What was crossed out is part of the document.
\newcommand{\struck}[1]{\sout{#1}}

% The transcriber's note, on the state of the page or the mathematics.
\newcommand{\note}[1]{\par\smallskip{\small\color{gray!60!black}\textit{[#1]}}\par\smallskip}

% A marginal note of hers, distinct from the body of the text.
\newcommand{\marginal}[1]{\par\smallskip\noindent\hspace{1em}%
  {\small\color{gray!60!black}#1}\par\smallskip}

% --- Title ------------------------------------------------------------------
% The title block is French, like the documents themselves.

\AddToHook{shipout/background}{%
  \ifx\grothWatermark\empty\else
    \begin{tikzpicture}[remember picture, overlay]
      \node[rotate=52, scale=3.4, align=center, text=gray!18,
            font=\sffamily\bfseries]
        at (current page.center) {\grothWatermark};
    \end{tikzpicture}%
  \fi}

\AtBeginDocument{%
  \begin{center}
    {\large\bfseries Manuscrits de mathématiciens (Gallica) ---
      \ifx\grothShelfmark\empty\grothFolder\else\grothShelfmark\fi}\\[2pt]
    {\itshape\grothFoldertitle}\\[4pt]
    {\small vues \grothFirst--\grothLast
      \ifx\grothBatch\empty\else\ (lot \grothBatch)\fi}\\[2pt]
    \ifx\grothDating\empty\else
      {\small\color{gray!60!black}Datation du catalogue : \grothDating}\\[2pt]
    \fi
    \ifx\grothWatermark\empty\else
      {\small\bfseries\color{red!45!gray}\grothWatermark}\\[2pt]
    \fi
    {\footnotesize\color{gray!60!black}Transcription établie d'après les images IIIF de Gallica.
     Source gallica.bnf.fr / Bibliothèque nationale de France.\\
     Les lectures incertaines sont soulignées, les illisibles notées [\,\ldots\,],
     les ajouts entre crochets ; \textsf{v.} numérote les vues de Gallica,
     \textsf{f.} la foliotation au crayon de la BnF.}
  \end{center}
  \vspace{1em}}

\endinput


\folder{nal-2338}
\shelfmark{NAL 2338}
\batch{3}
\pages{41}{60}
\dating{}
\watermark{Lecture automatique\\non vérifiée}
\foldertitle{Papiers divers de Mersenne.}

\begin{document}

\page{41}
\note{verso d'un feuillet ; la page continue, au milieu d'une phrase ouverte par une parenthèse, une lettre latine commencée avant cette vue. Aucun numéro écrit sur cette page : en haut à gauche, un « 20 » vu à l'envers par transparence, avec quelques mots de l'intitulé du recto (« \ldots{} Roberuallio \ldots{} »), qui ne se lisent pas de ce côté. Au-dessus de la première ligne, un grand trait de plume en crochet, plus foncé que le reste de la transparence : peut-être l'initiale ornée du recto vue au travers du papier ; ce n'est pas un mot.}

intentum meum) ipsius solidum demonstrare. Ipse vero dicis V. C\uncertain{l}: tandem habere Methodum pro reperiendo centro gr. plani ex data solidi planiq; mensura. Propositiones conuersae sunt. Inuersio autem huius modi facillima est; \& si in alterum ex nobis suspicio aliqua ferri debeat, certe in me non cadet; nam multo ante quam de hoc verbum faceres cum nostris Italis, sed \& cum Gallis, non solum enunciationem, sed etiam demonstrationem ipsam, orantibus vobis, ego misi in Galliam. Illud etiam pro me stare videtur argumentum, quod numquam ne verbum quidem fecisti de Centro gr. Cycloidis cum interea tantopere, \& quidem merito, gloriareis de omnibus alijs, Quadratura, Tangentibus, solidis, etiam de eo circa axem, quod tantum sperare dicebas. Verisimile non est cum reliqua omnia proponeres, \& in lucem edere velle promitteres quadraturas, tangentes, \& solida, de vnico Centro grauitatis siluisse, si illud tantum sperauisses quod quidem Problema meo iudicio nulli reliquorum post habendum videtur. sed de his, si opus fuerit multo plura dicemus suo tempore. Nemo tam facile suam laudem vnicuiq; tribuet quam ego dummodo tamen ignorantia, vel credulitate non decipiar. Methodum pro tangentib\textsuperscript{s} ex doctrina motus ego reperi pluribus abhinc annis, nulla ab alijs habita luce, vel auxilio: Cum amicis contuli; et in multis figuris propagaui. Postea incidi in demonstrationes Trochoidis et vtrumq; vulgaui inter amicos, antequam in meis libellis ederem. Ex improuiso, quando nil tale sperabam, nuncius horribilis ex vobis affertur haec omnia ante me vos etiam inuenisse. Si verum hoc est, certe pro meis illa amplius non essent habenda (quamquam fortasse nullus mortalium
\note{« V. C\uncertain{l} » : la troisième lettre, bouclée, peut être un l ou un s ; lu comme l'abréviation de « Vir Clarissime », que la lettre emploie en toutes lettres à la vue suivante. « gloriareis » est écrit ainsi. « promitteres » est écrit en deux morceaux, « promit teres ». « planiq; », « vnicuiq; », « vtrumq; » : q suivi du signe en forme de ʒ qui abrège -que. « gr. » abrège « grauitatis », écrit en entier plus bas. « inter » porte un petit trait suscrit. Les longs traits à la fin de plusieurs lignes sont des traits de remplissage, non des ratures.}

\page{42}
\note{recto ; la phrase de la vue 41 continue sans rupture. En haut à droite, à l'encre, « 21 » (voir l'en-tête : numéro de feuillet, non transcrit comme folio). Au bord gauche, de petites marques d'encre et des mots d'une autre page vus par transparence.}

ad haec vmquam descenderet). Vide Vir Clarissime quam ingenue ego agam cedendo, etiam ea quae iure atque mea sunt, ac vestra, cum vterq; proprio Marte adinuenerit, abstracta, (si qua intercesserit) modici temporis differentia. Sed incredibile est quanta iniuria afficiar dum video mihi praeripi ea, quae mea esse deberent sine controuersia Inuoco hominum fidem. Ecce verba ipsissima \uncertain{vestri} Clar\textsuperscript{mi} Mersenni in Epistola ad me data postquam enunciationem tantum Centri gr: Cycloidis sine demonstratione ad vos miseram dicebam enim secare axem in ratione 7 ad 5.
\note{« \uncertain{vestri} » : la dernière lettre porte un accent, et se lit aussi bien « vestrè ». « Inuoco » : l'o central est empâté.}

Dubitat Noster Roberuallus an Mechanicè tantum centra gr: Cycloidis, et semicycloidis inueneris, quae Geometrice falsa suspicatur docebis num istius rei demonstrationem habeas.
\note{ce paragraphe, en retrait, est la citation annoncée de la lettre de Mersenne ; aucun signe de citation ne le délimite.}

Quare ergo V. Clar\textsuperscript{me} dubitabas et Geometrice falsum suspicabaris quod ipse sciebas? Non dubito ego an vestra plana in infinitum abeuntia figurae genitrici aequalia sint, nam demonstrationem habeo. At ego non docui num demonstrationem illam haberem, sed protinus (quod quispiam praeter me non fecisset) arreptam demonstrationem in Galliam misi, simulque methodum vere meam, cum demonstratione pro inueniendo alterutro siue centro, siue solido in omnibus figuris, ex altero dato cum quadratura: quam methodum si apud vos ante habebatis, certe de veritate illa geometrica non erat quod dubitaretis, nam solidum circa basim tunc temporis vos habuisse dicetis, sed missa faciamus haec: Habeo enim et aliam methodum, quae vnica enunciatione determinat reperitq; Centrum grauitatis linearum superficierum ex reuolutione natarum, planorum, corporumq; omnium, dummodo axem, siue diametrum habeant. Vulgata
\note{« abeuntia » est récrit sur un premier tracé. Le point d'interrogation après « sciebas » est net.}

\page{43}
\note{verso du feuillet 21 ; « Vulgata » à la fin de la vue 42 se construit avec « est haec » qui ouvre cette page. En haut à gauche, « 21 » vu à l'envers par transparence ; plus bas, dans la marge de gauche, un trait de crayon ou d'encre pâle en forme de grand « L » bouclé, qui n'est pas un mot lisible.}

est haec apud amicos Italos; Oro vos ne inter vestra hanc etiam habeatis, nam hoc esset tollere penitus, omne litterarum scientiarumq; commercium. De libellis meis si tibi uir Clar\textsuperscript{me} placuerunt ut ipse affirmas (excepto tamen de Motu et Trochoide) maxima\add{s} ago gratias. In hoc autem plurimum \uncertain{discrepamus}, nam mihi ex ipsis nullus placet. Meliora non edidi, quia meliora habere non credebam: talia edidi, quasi coactus. Videbar enim quodamodo ad hoc teneri ex titulo officij mei Vtinam commercium hoc cum Clar: Galliarum Geometris, accidisset mihi ante aliquot annos, nam certe vel meliora edidissem vel nulla si reliqua bona essent leuis esset iactura libellum de Motu, et quidem totum ex meis opusculis aboleri.
\note{« inter vestra » est écrit d'un seul tenant, « intevestra ». « maxima » touche le bord du feuillet : la dernière lettre, que la construction demande, n'est pas sur la page. « \uncertain{discrepamus} » : l'avant-dernière voyelle a la forme d'un a (« discrepamas »). « quodamodo » est écrit ainsi.}

Argumenta quae producis V Clar\textsuperscript{me} contra fundamenta illius doctrinae fortia sunt, sed tamen non sunt noua. Extant huiusmodi obiectiones, immo etiam, \& multo fortiores apud ipsum Galileum in libro Proiectorum, sed latitant in illis Dialogis, quos ipse inter latinas Propositiones Italice interserit. Respondere possent eadem, quae respondet Galileus, et non nihil etiam ex meo; sed non aestimo tanti illum meum libellum, vt hanc molestiam suscipere tecumq; eandem partim V Clar\textsuperscript{me} aequum censeam Ad vltimum nemo hanc defensionem mihi eripiet, si dixero libellum esse ex maiori parte Geometricum. Archimedes supposuit olim proiecta non p. parabolas, sed p. lineas spirales suas procedere; et huius suppositionis gratia totum libellum composuit. At suppositionem esse falsam non multo post comperit Quid faceret? librum ne illum damnare penitus debebat? quin potius eadem omnia legi sineret sublata omni Proiectorum memoriä, additoq; vocabulo puncti cujusdam, quod non naturali
\note{« possent » est écrit au-dessus de la ligne, entre « Respondere » et « eadem », avec un signe d'insertion. « p. » rend deux fois le p barré de la copie (per). « partim » est lu avec la finale peu nette ; « censeam » a le c empâté. « memoriä » porte un tréma sur l'a, noté tel quel.}

\page{44}
\note{recto ; la phrase de la vue 43 continue sans rupture. En haut à droite, à l'encre, « 22 » ; au-dessus, sur le bord d'un feuillet qui dépasse, le haut d'un « 23 » (le feuillet suivant, vue 46). La moitié inférieure de la page est blanche.}

sed ficta quadam motuum lege moueatur. Proprium habet hoc Geometria quod nullis alienis auxilijs suffulta, sibi ipsi tantum innixa numquam ruere potest, sed mole sua stat, licet reliqua, omnia, quae propria culpa ipsi male adhaerebant, dilibantur. Reiciamus ergo quicquid Physicum est in eo libello nempe vocabula proiectorum, grauium, ballistarum et. c. et Geometricum maneat, hoc est propositiones abstractae. Caeterum Tabulae instrumenta, et si quid aliud hujusmodi est, non sint pro mensurandis proiectorum iactibus, sed pro determinandis quibusdam lineis geometricis, alijsq; ad geometricas parabolas pertinentibus. librum de sargues numquam vidi, sed quomodo probet ea quae proponit nescio. Videat ipse an liceat vnicuiq; talia scribere, \& fortasse etiam pari ratione probare, Videas autem ipse V Clar\textsuperscript{me}: quam facile sit hoc aeuo scriptorum feracissimo in aliorum reperta incidere, et pro suis innocenter aliena circumferre. Ipse proponis tamquam nouum quod Quadratrix transeat p. Centrum grauitatis semiperiphariae et. c. At multis abhinc annis hoc ipsum demonstratum fuit, et editum ab alijs multo vniuersalius.
\note{« dilibantur » est écrit ainsi. « de sargues » : le nom est écrit en deux mots, sans majuscule. « et. c. » est la forme de la copie pour « etc. ». « semiperiphariae » : la finale est la ligature æ de la copie.}

Hoc ipsum ostenderam ego quoque, maxima, et breuitate, \& simplicitate, antequam scirem quemdam Guldinum Iesuitam haec edidisse. Monitus enim fui nuperrime a Clar\textsuperscript{mo} Viro Caua\uncertain{serio} me in idem inuentum incidisse atq; adeo perdidi contemplationem qua nullam habeo elegantiorem, aut magis absolutam. Admiratus sum subtilissimum Theorema de spatijs in infinitam longitudinem abeuntibus. Egregium inuentum, et quod nugis meis vel ipse preferam. At non dico inter mea illud habuisse cum ad infinitam illam distantiam (exceptis hyperbolis
\note{le nom propre est coupé en fin de ligne : « Caua » à la fin de l'une, puis, en tête de la suivante, un mot dont la première lettre a la forme de l's long de la copie, « \uncertain{serio} ». La lecture « Cauallerio » n'est pas celle de la page ; elle n'est pas substituée. Un petit « z » ou « 2 » isolé à l'encre, au bord droit, à hauteur de « Caua » : peut-être un signe de renvoi ; sens non établi.}

\page{45}
\note{verso du feuillet 22 ; la parenthèse ouverte au bas de la vue 44 se ferme ici. En haut à gauche, « 22 » vu à l'envers par transparence ; le haut de la page laisse voir, par transparence, les lignes du recto.}

meis) non \uncertain{respexerim}. Circa illius demonstrationem nihil prorsus moratus sum; nam vix lecta propositione statim animaduerti ipsi conuenire demonstrationem, qua iam pridem vsus fueram in dimensione infinitarum parabolarum. Vulgata iam est apud amicos Italos, ipsamq; hic obiter adiungam orsus a Tangentibus, quando quidem tangentes in ea adhibentur: quae omnia (licet ineptiae sint) fidei Vestrae commendo.
\note{« \uncertain{respexerim} » : la troisième lettre a la forme d'un x (« rexpexerim »).}

Lemma I.

Esto parabola qualibet ABC ad cuius punctum A danda sit tangens. Applicetur AD ad diametrum ED, fiatq; vt exponens applicatarum, qui exemp: gr: sit 5 ad exponentem diametralium qui exemp: gr: sit 3 ita recta ED ad DC dico iunctam AE non conuenire cum parabola ad aliud punctum praeter A. Nam si potest conueniat in B, \& applicata B.I, fiat vt ED ad DC, ita EI, ad IO sitq; latus rectum CL, cui aequidistet OP.
\note{« Esto » s'ouvre sur une grande initiale ornée. Aucune figure n'est tracée sur la page, bien que le texte en suppose une (points A, B, C, D, E, I, O, P, L). « 5 » est écrit d'une forme proche d'un ς.}

Erit iam dignitas B.I, hoc est in nostro casu cuboquad\textsuperscript{m} B.I aequale ductui dignitatum. IC in CL, nempe ductui Cubi IC in quadratum CL ob latus rectum. At idem cuboquad\textsuperscript{m} BI aequatur ductui cubi IO in quadratum OP (vt mox patebit) ergo ductus IC in CL, et IO in OP sunt aequales. Quod est impossibile, vt demonstrabimus inferius.
\note{« cuboquadratum » (le carré-cube, cinquième puissance) répond à l'exposant 5 des appliquées posé plus haut, et « cubi IC in quadratum CL » aux exposants 3 et 2 : les deux exemples concordent. Le point entre « dignitatum » et « IC » est sur la page.}

Quod promisimus primo ostendetur sic. Cuboquadratum AD, ad cuboquad\textsuperscript{m} BI rationem habet compositam ex ratione cubi

\page{46}
\note{recto ; la phrase de la vue 45 continue sans rupture. En haut à droite, à l'encre, « 23 ». La moitié inférieure de la page est blanche.}

AD, ad BI, siue cubi DE ad EI; siue cubi DC ad IO. Et ex ratione quadrati AD ad BI siue quad\textsuperscript{ti} DE ad EI, vel quadrati CE ad EO, hoc est quadrati CL. ad OP. Ergo cuboquad\textsuperscript{m} AD. ad BI erit vt ductus cubi DC in quadratum CL ad ductum cubi IO in quadratum OP. Sed antecedentia sunt aequalia ob latus rectum, ergo \& consequentia aequalia erunt quod promisimus primo.
\note{« quadrati » est écrit en deux morceaux, « qua diati ».}

Secundum vero ostendimus sic. Recta IO ad OE est vt exponens diametralium ad differentiam exponentium, nempe vt 3 ad 2 Quare ductus cubi IO in quadratum OE erit maximus. (quod demonstrauimus iam pridem vniuersaliter, geometrice sine vlla ope Algebrae) sed vt ductus IO in OE ad ductum IC in CE ita ductus IO in OP ad IC in CL. Quare et IO in OP maximus erit.
\note{3 ad 2 : l'exposant des diamétrales (3) à la différence des exposants (5 $-$ 3), selon les exemples posés au lemme I (vue 45). Le texte de la page écrit « ductus IO in OE » et « IO in OP » là où la phrase précédente parle du cube de IO et du carré de OE ; il est transcrit tel qu'il est.}

Haec una est ex quinque Methodis quas habemus pro tangentibus.

Lemma II.

Omnia parallelogramma puta AB, CD ad inscriptas sibi parabolas (dumodo parabolae sint eiusdem speciei) eandem rationem habent. Ponimus tantum figuram, \& omittimus demonstrationem, nam facillima est, praecipue vero p. doctrinam Indiuisibilium, addita tertia parabola in parallelogr\textsuperscript{m} B.D, et ductis vbicunq; duabus tantum rectis \struck{E, F, G H.} e, f, g, h.
\note{l'énoncé du lemme II est écrit sur la moitié gauche de la page seulement, la droite restant blanche, comme pour une figure ; « Ponimus tantum figuram » l'annonce, mais aucune figure n'est tracée. Les majuscules « E, F, G H » sont barrées et suivies des minuscules « e, f, g, h ». « dumodo » est écrit ainsi.}

Iam quadratura parabolarum talis est. Esto parabola qualibet ABC, cuius diameter BC, basis AC parallelogrammum vero

\page{47}
\note{verso du feuillet 23 ; la phrase de la vue 46 continue sans rupture. La moitié gauche de la page laisse voir, à l'envers, les lignes du recto. Aucune figure n'est tracée. Trois passages sont soulignés d'un trait d'encre continu, sans qu'on puisse dire si le trait est du copiste ou d'un lecteur : « linea Roberualliana » jusqu'à « hoc fateri », « eandem demonstrationem » jusqu'à « nunc mitto », et « Estq; haec demonstratio » jusqu'à « ad te Mitto ».}

circumscriptum sit FC. Esto Tangens. AD; concipiaturq; in parallelogrammo FB.DE alia parabola eiusdem speciei, siue quod idem est linea Roberualliana (noua enim illa linea quam ipse excogitasti Vir Clar\textsuperscript{me} si ortum ducat ex aliqua parabolarum semper parabola euenit eiusdem speciei, quod ego nouum esse scio, licet fortasse turpe videatur hoc fateri) \& erit parabola ABC aequalis trilineo ABE, vt ego demonstraui ellapso iam biennio p. eandem demonstrationem quam ad tuas lineas applicatam in folio separato ad te nunc mitto.
\note{« Esto Tangens. AD; » : le point est bien entre « Tangens » et « AD ». « excogitasti » est coupé en fin de ligne (« ex / cogitasti »). « ellapso » est écrit ainsi. Le « folium separatum » annoncé ici n'est pas dans les vues 41–60.}

Iam vt exponens applicatarum ad exponentem diametralium ita est recta DC ad CB ob tangentem p. lemma p\textsuperscript{m}. Et ita est parallelogrammum EC ad FC. Sed ita etiam sunt ambae simul parabolae EBD, ABC ad parabolam ABC, ob 2\textsuperscript{m} lemm. Ergo p. 19 Quinti reliquum EBA siue parabola ABC ipsi aequalis, erit ad reliquum FBA vt erat totum ad totum hoc est vt exponens applicatar\textsuperscript{m} ad exponentem diamet\textsuperscript{m}. Propterea conuertendo, componendoq; patet parallelogrammum quodlibet ad suam parabolam esse vt aggregatum exponentium ad exponentem applicatarum.
\note{« lemma p\textsuperscript{m} », « 2\textsuperscript{m} lemm. » : le lemme I (vue 45) et le lemme II (vue 46). « 19 Quinti » : la proposition 19 du livre V d'Euclide, selon toute apparence.}

Sed \& alia ratione concludimus maiori breuitate quadraturam omnium parabolarum. Estq; haec demonstratio veluti Corollarium demonstrationis quam in folio separato ad te Mitto. Videbis vir Clar\textsuperscript{me} ex eodem inuento ex quo gloriaris te quadraturam vnius parabolae quadraticae deduxisse, nos quadraturam omnium nullo plane negotio deriuare

\page{48}
\note{recto ; la phrase de la vue 47 continue sans rupture. En haut à droite, à l'encre, « 24 ». Comme au lemme II, la seconde démonstration est écrite sur la moitié gauche de la page, la droite restant blanche, sans figure. Passages soulignés d'un trait d'encre : « ignotum sit » jusqu'à « cum genitrice », « linea noua Roberualliana », « non tamen patiar mihi illas eripi ».}

p. indiuisibilia more nostro Quodq; elegantius est demonstratio concludit quamquam ignotum sit cuius nam speciei sit noua illa linea tua, quam diximus supra parabolam esse eiusdem gradus cum genitrice.

Esto parabolarum aliqua BAFC cuius diameter AB, basis vero BC parallelogrammum circumscriptum BL, linea noua Roberualliana quaecunq; tandem sit esto ADE, \& sūpto quolibet puncto F in linea parabolica, sit applicata FI, Tangens F.H, et ipsa GFD diametro aequidistet: tum iuncta DH applicatis aequidistans et ipsa erit ob definitionem figurae.
\note{« sūpto » : « sumpto », avec le tilde de la copie sur l'u. Les lettres O et G, nommées ici et plus bas, ne sont définies sur la page par aucune figure.}

Iam DF recta ad FO erit vt HI, ad IA, siue ob Tangentem, vt exponens applicatarum ad exponentem diametralium. Et semper ita continget vbicumq; sumpta fuerit recta DF. Ergo omnes antecedentes simul, nempe trilineum EDAFC, siue parabola ABC ipsi aequalis, erunt ad omnes consequentes simul, nempe ad trilineum L.AC, vt exponens applicatarum ad exponentem diametralium. Quadratura ergo patet conuertendo, componendoq; vt supra.

Praedicta Methodi tum pro Quadraturis, tum pro Tangentibus sunt quas minimi prae caeteris ego facio (non tamen patiar mihi illas eripi) Supersunt adhuc pro Tangentibus quatuor, pro quadraturis vero adhuc tres Methodi, inter quas vna est praecipua, et

\page{49}
\note{verso du feuillet 24 ; la phrase de la vue 48 continue sans rupture. Aucun numéro écrit sur cette page. Un passage souligné : « et hoc Geometrice more veterum demonstratum est a nobis ».}

generalissima, quae communis est \& parabolarum, \& hyperbolarum tam quadraturis, quam etiam tangentibus, alijsq; Problematibus, quae nundum circumferre est animus.

Iam examinauimus figuras planas, solidasq; in infinitum recta via excurrentes. Alium modum habet Geometria in infinitum abeundi, \& quidem mensurabilem Problema circa lineam est, quae cum sit curua, et ex infinitis spiris conflata, rectae tamen lineae aequalis est; et hoc Geometrice more veterum demonstratum est a nobis. Figuras quasdam huius Problematis sine vlla definitione misimus in Galliam iam ante biennium. Clar\textsuperscript{mo} tamen Mersenno dum esset Romae communicauimus quasdam \uncertain{rationes} pro descriptione proxima earumdem spiralium, quarum quaedam Theoremata vna cum definitione ad vos nunc mitto. Videbatur aliquid deesse, cum reperta solidorum et spatiorum in infinitum abeuntium dimensione, nulla de lineis adhuc facta esset mentio. Linearum quidem genus non potest in infinitam distantiam abire quin \& quantitas infinita sit. At potest occultum habere caput et inter laberintos infinitarum reuolutionum adeo reconditum vt solis Geometriae oculis sit visibile. Quicquid est, \& an Problema dignum sit oculis tuis ipse videbis Vir. Clarissime.
\note{« nundum », « laberintos » sont écrits ainsi. « \uncertain{rationes} » : la première lettre, en tête de ligne, peut être un r ou un v de la copie (« vationes »). Le « Problema » est noté en marge gauche d'un trait horizontal, sans mot. Les « Theoremata » et la « definitio » des spirales annoncés ici ne sont pas dans les vues 41–60.}

Hyperbolarum Theoremata, quae mitto ad Ill\textsuperscript{m} De Ferma vt iudicium subeant, num cum parabolis saltem aliqua ex parte conferri possint, videre poteris. Si vnius hyperbolae Primariae Quadratura tam diu quaesita est, Nos pro vna
\note{ce paragraphe s'ouvre sur une grande initiale à cadeaux. « De Ferma » est écrit ainsi, sans t, à la fin d'une ligne dont rien ne manque au bord.}

\page{50}
\note{recto ; la phrase de la vue 49 continue sans rupture. En haut à droite, à l'encre, « 25 ». Seul le haut de la page est écrit ; le reste est blanc. Aucune signature.}

infinitas damus. Interea toto affectu me tibi commendo V. Clarissime.

Vale. Dat. Florentiae die 7 Julij ann. 1646
\note{le millésime est souligné d'un trait. La lettre finit ici, sans souscription ni signature ; son destinataire est Roberval, que le texte interpelle (« noua illa linea quam ipse excogitasti Vir Clar\textsuperscript{me} », vue 47 ; « linea tua », vue 48). Son auteur n'est pas nommé dans les vues 41–60.}

Factum etiam est vt vsque ad initiu\add{m} Augusti Mensis Litterae \struck{iam} \add{in}tabulatis apud me morarentur ob varia impedimenta mea.
\note{post-scriptum, ouvert par une grande initiale ornée. « in » est écrit au-dessus de « iam » barré. « initiu » porte une queue d'abréviation pour la finale. Les dernières lettres, « a mea », sont repassées d'une encre plus noire. Le post-scriptum, qui dit la lettre retenue jusqu'au début d'août, est postérieur à la date du 7 juillet qui le précède ; la page ne dit pas de combien.}

\page{53}
\note{recto ; une nouvelle lettre commence ici, de la même main de copiste que la précédente. En haut à droite, à l'encre, « 26 ». Au milieu de la page, le cachet rond « BIBLIOTHÈQUE NATIONALE MSS » avec « R.F. ». Le bas de la page est blanc sous la dernière ligne.}

Ill\textsuperscript{mo} et Doctissimo Viro P. de Carcauy\\
Euang. Torricellius Sal.
\note{l'intitulé donne à la fois le destinataire et l'expéditeur ; c'est le seul nom d'auteur que portent les vues 41–60.}

Circa Problema numericum Ill\textsuperscript{ri} Senatoris De Fermat nihil moratus sum: totus enim, alienus à studijs omnibus fui integro hoc anno et fortasse etiam in sequentibus ero, cum alia mihi vitae ratio ineunda sit Dubitaui etiam ne problemata ista numerica, quae \add{c}ommunem, et vulgatam Algebrae methodum fortasse excedunt, difficilis ad modum solutionis essent; praesertim si quis illa tantum inquirat data opera, quamquam postea se se offerant, processu temporis quando ea disciplina colitur ex instituto, et assidua contemplatione euoluitur. Praeterea non tam plausibile mihi videbatur inuentum illud. Omnes potestates quarum exponens \&c. Si vnitate augeantur numeros primos fieri: illudque. Triangulum rectangulum in numeris reperire cujus bina latera quadratum efficiant, siue alia simili conditione propositum, quod non memini, vt operae pretium ducerem ingeniolum meum propriae gloriae adeo indigum circa alienam diutius torquere.
\note{« Circa » s'ouvre sur une grande initiale C ornée, qui occupe la marge sur cinq lignes. « integro » est d'une encre plus pâle que le reste de la ligne. « \add{c}ommunem » : la ligne commence par « ommunem », sans c visible. « Ill\textsuperscript{ri} » : l'abréviation est écrite en exposant, lue ri.}

At non hujusmodi visum, est problema tam vastum, tam multiplex de infinitis parabolis, quod ego inuentum primo existimabam a Cl. Roberuallio proficisci: deinde audiui ab Ill\textsuperscript{mo} D. de Fermat repertum postremo comperio Cl. Caualerium hoc inuentum sibi vindicare, cum ipse multis abhinc annis illud contemplatus sit, quamquam frustra, deinde anno 1639 in publicam lucem emiserit in quodam suo libello proposuerit que: reperiendum per idem fere tempus Clar\textsuperscript{mo} Viro I F. niceroni dum Bononiam pertransiret communicauit vt in Galliam ferret, quamquam ipse Caualerius absolutam solutionem non dum haberet. Sed nihil hoc ad nos attinet. hoc vnum sciat \uncertain{verum} Ill\textsuperscript{mus} D. De Fermat me demonstrationes
\note{« I F. niceroni » : « I F » en fin de ligne, suivi d'un trait d'union, puis « niceroni » en tête de la suivante. « \uncertain{verum} » : le mot se lit aussi « venim », l'r de la copie ressemblant à un n. « proposuerit que: » est écrit ainsi, l'enclitique détaché.}

\page{54}
\note{verso du feuillet 26 ; la phrase de la vue 53 continue sans rupture. En haut, l'intitulé du recto transparaît à l'envers. À droite, au-delà de la marge, le haut d'un « Ill » appartient à la page suivante, visible sous ce feuillet (vue 55). Sur quatre lignes, au milieu de la moitié droite (lignes 6 à 9), l'encre a pâli jusqu'à presque disparaître ; les mots perdus ont été cherchés à plusieurs contrastes sans être retrouvés en entier.}

omnes circa praedictas parabolas (cujuscunque sint) reperisse vniuersalissimas, licet nescio quo pacto definitio exciderit non adeo vniuersalis. Quod ad tangentes attinet quinque methodos habeo penitus inter se diuersas Quadraturas etiam totidem, ex quibus iam \uncertain{duas} quas prae caeteris minime facio vulgaui inter amicos, et ad Clar. Roberuall\uncertain{ium} \ill{}to, fortasse ad subeundam eandem fortunam cum meo \uncertain{centro} \uncertain{g}rauitatis Cycloidis. Circa centrum gr. vestri Geometra\uncertain{e} p\ill{}ant mira se habere, et vniuersalissima. Ego vero \uncertain{quid} habeam praeter solitas nugas meas. Misi iam ante biennium ad Cl. Roberuallium demonstrationem methodi, meae pro reperiendo siue centro, siue solido alicujus plani ex altero tantum dato vna cum quadratura. Exponam hic enunciationem alterius cujusdam Propositiunculae, non vt demonstratio inueniatur nam facillima est, et iam vulgata inter amicos Italos, sed vt magis elucescant aliorum amplissimae contemplationes.
\note{« \uncertain{duas} » : le mot porte une lettre de trop, lue comme un s long (« dufas »). Dans la zone pâlie : après « Roberuall » on devine « ium », puis un mot de deux ou trois lettres et la fin d'un verbe en « -to » (« mitto » ?), non restitués ; « \uncertain{centro} » et le g de « grauitatis » sont à peine visibles ; après « Geometra » on distingue un a final, puis un p initial et la fin « -ant » d'un verbe de quatre à six lettres illisibles ; « \uncertain{quid} » se réduit à un q et à des traces. « Cl. » rend un C suivi d'un l bouclé et d'un point.}

In quacunque figura siue plana, siue solida dumodo axem, vel diametrum habeat, centrum grauitatis axem, siue diametrum secat vnica semper hac lege: Nempe vt pars axis vel diametri versus verticem figurae sit ad reliquam quemadmodum sunt omnes \&c; ad omnes \&c.

Haec est regula ex qua centra grauitatis exprimo, cum habeam methodum non adeo difficilem pro inuenienda ratione quam habent pred: Omnes \&c. ad omnes \&c. Immo sub eodem Theoremate compraehenduntur centra gr. linearum, et superficierum paucis tantum mutatis in Enunciatione.
\note{« \&c » : la copie laisse l'énoncé incomplet (« omnes \&c. ad omnes \&c. ») ; les termes omis ne sont pas sur la page. « pred: » abrège « praedictam ».}

Sed quid est cur tantopere petatis iudicium meum de

\page{55}
\note{recto ; la phrase de la vue 54 continue sans rupture. La vue montre à gauche le bord du feuillet 26v (fins de lignes de la vue 54, non répétées ici) et, à droite, le feuillet suivant, numéroté à l'encre « 27 » en haut à droite. La lettre finit sur cette page ; l'adresse est écrite au pied, à gauche, plus bas que la souscription.}

Aristarchi libello? idem postulauit Cl. Mersennus dum esset Romae Amici mei existimant libellum plane diuinum et ab Auctore diuino compositum. Ego censeo libellum sub Aristarchi nomine editum, conscriptum fuisse nostra haec aetate. Quod attinet ad doctrinam, omnia quidem optima credo cum à doctissimis viris probentur, attamen, et mihi, \& quibusdam amicis quam plurima non placent ob ingenij nostri imbecillitatem. Sed quaeso ne et rationes postuletis, quemadmodum fecit ipse. Cl. Mers. cur ego libellum nuper conscriptum censeam, siue cur in eo multa displiceant. Ridiculum sane esset me multa potiora, et ad me spectantia consulto negligentem, circa negotium quod ad me minime attinet excruciari. Oro D. V. vt inuentum meum de infinitis hyperbolis, et si placet etiam de spiralibus statim innotescat non solum Ill\textsuperscript{mo} de Fermat, sed etiam alijs Geometris. Quando enim ego misi iam ante biennium demonstrationem de Centro gr. Cycloidis cum demonstratione Methodi pro reperiendo siue centro siue solido alicujus plani, memini me orauisse Cl. Mersennum vt vtramque demonstrationem cum multis statim conferret; quod si ille fecisset certe nunc mihi mea non eriperentur, quae alij mihi debent, nam primus inueni, imo solus inueni. Vale Vir Ill\textsuperscript{me}, et me inutilem quidem, sed obsequentissimum famulum vt cepisti ama Euangelista\textsuperscript{m} Torricellium. D. flor. Die \uncertain{8} Julij an. 1646.
\note{« Oro D. V. » : « Dominationem Vestram », selon l'usage. « Euangelista\textsuperscript{m} » : le m final est suscrit. Le quantième du jour est un chiffre fait de deux petites boucles côte à côte, qui se lit « co » ; lu 8 avec doute. La lettre précédente, à Roberval, est datée du 7 juillet 1646 (vue 50).}

Ill\textsuperscript{mo} Viro P. de Carcauy

\page{57}
\note{recto ; une troisième pièce commence ici, d'une autre main que les deux lettres précédentes : écriture plus petite, plus liée, avec les accents de l'époque sur les voyelles (à, â, ò, ù). En haut à droite, à l'encre, « 28 ». Au-dessus du titre latin, deux lignes en français, barrées, d'une encre et d'une main qui semblent être celles de la note du pied de page. Au milieu des lignes 1 à 7, l'encre a pâli sur une large zone, avec une tache sombre ; les mots de cette zone ont été cherchés à plusieurs contrastes et sont en grande partie perdus. Le cachet rond « BIBLIOTHÈQUE NATIONALE MSS » avec « R.F. » est posé sur le texte au milieu de la page.}

\struck{Lettre de M\textsuperscript{r} De Roberval escrite au S\textsuperscript{re} Toricelli}

Epistola Aegidii Personerii de Roberval\\
ad \struck{R. P.} Euangelistam Torricellium
\note{« \struck{R. P.} » : deux majuscules barrées entre « ad » et « Euangelistam ». Dans la ligne française barrée, « Toricelli » et « escrite » sont lus sous la rature.}

Vir clarissime

Si me vnum respicerem: si nullâ \ill{} \ill{}m me; si nullâ ipsius, quam prae caeteris \ill{} \ill{}tâ ratione, internâ animi tranquillitate conquiescerem: non me\ill{} \ill{} profectò, quòd vos deûm atque hominum fidem inuocetis, quòd celeber\ill{}m hominum testimonium in me adducere conemini; quòd denique null\ill{} \ill{}n moueatis lapidem, ad hoc vt ego meorum ipsius operum plagiarius ho\ill{} quippe qui planè mihi conscius sum, ex ijs quae ad vos scripsi, nihil non ve\ill{}sse. Sed fateor ingenuè; longè absum à praestanti illo vitae philosophicae statu; tantámque beatitudinem si optare nobis licet, non etiam sperare statim licet. Ego enim inter multos natus, inter multos educatus, cum multis viuere atque conuersari assuetus, cum multis etiam necessitudines \uncertain{contraxi}; ita vt rebus externis non moueri huc vsque nondum didicerim. Itaque admonet nos existimatio nostra, quam tueri, quámque, si quo id labore liceat aut impendio, promouere tenemur: postulant amici, collegae, Mathematici Galliarum praestantissimi, quibus omnia me debere fateor: cogit ipsa cui totum me dicaui veritas. ne tam grauem vestram accusationem prorsùs negligam; praesertim quam nullius negotij fuerit refellere; cùm praeter rationes nostras, quae per se sufficiunt, ijsdem ambo testibus vtamur. Erit etiam quod de vobis expostulem; et vt spero non injuriâ; qui cùm festucam in nostris oculis quaeratis, trabem in vestris non animaduertatis. Nolim tamen ob id, tolli inter nos literarum commercium; quòd vos nimium rigidè, meo quidem judicio, quasi aliquid nobis timendum minati estis: quin potiùs optarim tales iras, suauissimi commercij redintegrationem esse. Quòd si inter nos, per nosipsos conuenire non potest, judicent amici, nos judicio ipsorum stare promittamus. Ad rem venio.
\note{la zone pâlie emporte : la fin de la ligne 1 après « si nullâ », sauf « \ldots{}m me; » ; après « prae caeteris » un mot commençant par u et le début d'un mot finissant en « -tâ » ; la fin de « non me\ldots{} » et le mot qui suit ; le milieu de « celeber\ldots{}m » (un génitif pluriel, selon la construction) ; après « null- » une ou deux syllabes, puis la fin « -n » d'un mot bref avant « moueatis lapidem » ; après « plagiarius ho- » la fin du mot et un ou deux mots ; entre « ve- » et « -sse » environ un mot. Aucune de ces lacunes n'est comblée. « \uncertain{contraxi} » : la forme se lit ainsi, bien que la construction attende un infinitif. « conuersari » est écrit avec deux s longs très ouverts.}

(De propositione Rotae atque Trochoidum illius, primùm audiui Parisijs anno 1628. (eo enim demùm anno, ab expeditione Rupellana reuersus, statui in maxima illa atque omni studiorum genere excultissima vrbe, firmas sedes stabilire; cùm anteà vagus, incertis sedibus, diuersis in regni Gallici partibus degissem) asseruítque qui proponebat celeberrimus vir pater Mersennus, talem quaestionem per multos jam annos à pluribus tentatam, eousque insolutam permansisse: cui ego respondi, hoc ei commune esse cum multis alijs vetustissimis nobilissimisque propositionibus: neque ideò quicquam in illa magis quàm in his mirandum videri, si vnâ cum illis Solutione careret. Ac tunc ipse,
\note{la parenthèse qui ouvre ce paragraphe, devant « De propositione », n'est pas fermée sur cette page ; elle est transcrite telle qu'elle est. Le reste de la page, sous « Ac tunc ipse, », est blanc jusqu'à la note du pied.}

\marginal{Vous prendrez garde de faire de caracteres Italiques ce qu'il y a icy de gros caracteres}
\note{au pied de la page, à droite, en français, d'une autre main et d'une encre plus pâle que le texte : une consigne sur les caractères, adressée, semble-t-il, à un imprimeur ou à un copiste ; à quels passages elle renvoie n'est pas dit. Ce n'est pas une note de l'auteur de la lettre ; elle est donnée en marge faute de mieux. « Vous » est lu avec une initiale en forme de J.}

\page{58}
\note{verso du feuillet 28 ; la phrase de la vue 57 (« Ac tunc ipse, ») continue ici. Aucun numéro ; en haut, l'intitulé du recto transparaît à l'envers. Au début des lignes 4 à 7 et au milieu des lignes 6 à 9, l'encre est pâlie, mais se lit encore.}

cùm difficillimam existimarem, certè supra vires meas, intactam ita dimisi vt per sex annos de illa ne quidem somniarim. Atque vt verum fatear, ego tunc annum agens 27\textsuperscript{um} etiamsi continuo decennij anteacti exercitio, discendo docendóque, atque agendo, in rebus Mathematicis, in primis verò in Analyticis, quibus etiamnum maximè delector, non mediocriter profecissem; tamen, neque eum adhuc habitum mihi comparaueram; neque eas ingenij vires susceperam, quae ad eiusmodi quaestiones sufficerent. Intereà, cùm mecum ipse saepiùs cogitarem, quâ potissimùm ratione possem in suauissimae Matheseos adyta penetrare; statui diuinum Archimedem, quem ferè vnum inter antiquos Geometras suspicio, attentiùs considerare: ex qua consideratione sublimem illam et numquam satis laudatam infiniti doctrinam mihi comparaui: sic enim tunc vocabam eam quae à clarissimo Cauallerio vocatur doctrina indiuisibilium. Ridebis forsan; \uncertain{St.} Hic ergo Gallus, inquies, non solùm Trochoidum dimensionem ante nos, si Dijs placet: non solùm Parabolarum omnium: non solùm solidorum ad has et illas pertinentium: non solùm planorum ab helicibus cujuscumque gradus aut dignitatis compraehensorum: non solùm earumdem helicum secundùm longitudinem, cum praedictis Parabolis comparationem: non solùm curuarum omnium tangentes per motuum compositionem: non solùm doctrinam centrorum grauitatis inuenerit? sed et praestantissimi nostri Cauallerij indiuisibilia quoque? atque illa omnia nobis; haec illi, plagiarius ille impunè eripuerit? Verumtamen rideatis licet, et talia, aut ijs pejora de nobis putetis, aut vociferemini: ego Trochoides, parabolas, helices, tangentes, et centra ante vos; imò et multò plura non solùm inueni, sed et vulgaui. An vultis vt verum reticeam, quod partes nostras adjuuat: falsum autem proferam, quod nobis nociturum sit? Nos aetate aut tempore saltem, priores, aetatis aut temporis beneficia respuemus? et junioribus aut saltem tempore posterioribus, viui adhuc relinquemus? apage stultam illam in nosmetipsos injusticiam. Quòd si cuncta ego vnicâ epistolâ quam ad vos scripsi, non enumeraui, nihil mirum; illa enim aliunde satis prolixa extitit: nec id necessarium, aut operae pretium judicaui. Deinde etiam, quid de paucis aliquot propositionibus enumeratis gloriari attinet?

\begin{quote}
Pauperis est numerare pecus.
\end{quote}

Sed de vobis plura posteà; nunc de indiuisibilibus, quoniam illa ad rem faciunt, dicamus. Illa ergo, an ante nos clarissimus Cauallerius inuenerit, nescio: certè illud scio, me integro quinquennio antequam in lucem emiserit, eâ doctrinâ vsum fuisse in soluendis multis, ijsque planè arduis propositionibus. Attamen, absiste moueri: ego tanto viro, tantae ac tam sublimis doctrinae inuentionem non eripiam; nec possum; nec si possim, faciam. Ille prior vulgauit: ille, hoc jure,
\note{« 27\textsuperscript{um} » : « vigesimum septimum ». « \uncertain{St.} » : deux lettres et un point en fin de ligne, lues « St. » (l'interjection) ; « Et. » n'est pas exclu. Les points d'interrogation de cette main ont la forme d'un crochet sur un point. « Pauperis est numerare pecus » est centré sur sa ligne, comme une citation (Ovide, \emph{Métamorphoses}, XIII, 824, selon toute apparence ; la page ne le dit pas).}

\page{59}
\note{recto ; la phrase de la vue 58 (« ille, hoc jure, ») continue ici. En haut à droite, à l'encre, « 29 ». Même main qu'aux vues 57–58.}

suam fecit: ille, hoc jure, habeat atque possideat: ille tandem, hoc jure, inuentoris nomine gaudeat. Absit vt in posterum, quod nec priùs feci, in tali causa, intercessoris ridiculi prouinciam mihi suscipiam; praesertim cùm nequidem inter amicos quicquam vnquam de tali doctrina vulgauerim, quam neque publici juris facere, nisi post aliquot annos, juuenili quodam mei ipsius amore, decreueram. Quippe sperabam interim fore vt solutione difficiliorum quaestionum quas quotidie nullo negotio, tali instrumento adjutus, vulgabam, doctrinae famam facilè consequerer: neque sanè haec spes ex toto me fefellit: postquam enim ingenti ardore doctrinam ipsam excoluissem, eamdémque ad puncta, ad lineas, ad superficies, ad angulos, ad solida praecipuè; postremò etiam ad numeros extendissem, haud fuit difficile ea exequi propter quae amici laetarentur, inuidi disrumperentur. Exultabam ergo nimiùm juueniliter, ac tantò diligentiùs doctrinam ipsam reticebam; dignus planè in quem poeta dixerit

\begin{quote}
Nec ferre videt sua gaudia ventos.
\end{quote}

qui detectâ auri fodinâ ditissimâ; dum grana quaedam ex ea decerpta ostento, vt ex diuitibus ac beatis quidam habear; interim alius eamdem à se quoque detectam, palàm, plaudentibus omnibus, ostendit, ac publici juris facit: ita vt exinde periculum sit ne ridear, si à me quoque inuentam fuisse affirmauero. Est tamen inter clarissimi Cauallerij methodum et nostram, exigua quaedam differentia. Ille enim cujusuis superficiei indiuisibilia secundùm infinitas lineas; solidi autem indiuisibilia secundùm infinitas superficies considerat. vnde ex vulgaribus geometris plerique; sed et quidam ex superbis illis sciolis qui soli docti haberi volunt, quique si nihil aliud, certè hoc vnum satis habent vt in magnorum virorum opera insurgant, quòd à se minimè profecta esse inuideant, occasionem carpendi Cauallerij arripuerunt, tanquam si ille aut superficies ex lineis, aut solida ex superficiebus reuerà constare vellet. Quamquam autem illi coram eruditis nihil aliud lucrentur quàm ignorantiae, aut inuidiae titulum, tamen ijdem coram imperitis, suâ authoritate, de doctorum famâ non mediocriter detrahunt; nec ab ijs illaesus euasit Cauallerius. Nostra autem methodus, si non omnia, certè hoc cauet, ne eterogenea comparare videatur. Nos enim infinita nostra seu indiuisibilia sic consideramus. lineā quidem, tanquam si ex infinitis seu indefinitis numero lineis constet; superficiem, ex infinitis seu indefinitis numero superficiebus; solidum, ex solidis; angulum, ex angulis; numerum indefinitum, ex vnitatibus indefinitis. Imò plano-planum, ex plano-planis numero indefinitis componi concipimus; atque ita de altioribus: singula enim suas habent vtilitates.
\note{le vers centré est cité comme d'un poète que la lettre ne nomme pas. « superficiei » est écrit « superficiej ». « eterogenea » est écrit sans h. « lineā » : le trait sur l'a (lineam) est celui de la page.}

Dum autem speciem aliquam in sua infinita resoluimus, aequalitatem quandā, vel certè notam aliquam progressionem inter partium altitudines aut latitudines ferè semper obseruamus. Sed de hoc satis superque: nunc ad vos redeo. Cùm itaque ope indiuisibilium multa protulissem;
\note{« protulissem » : les lettres du milieu sont surchargées, deux s longs récrits sur d'autres lettres. La page finit sur ce point-virgule ; la phrase continue au-delà de la vue 60, selon toute apparence (voir la vue 60).}

\page{60}
\note{verso du feuillet 29 ; la phrase de la vue 59 (« Cùm itaque ope indiuisibilium multa protulissem; ») continue ici. Aucun numéro. Même main qu'aux vues 57–59.}

tandem anno 1634. celeberrimus P. Mersennus Trochoidem in memoriam reuocauit, non sine graui expostulatione, quasi propositionem haud quaquam ignobilem, de industria praeterirem difficultate illius perterritus: Ego sic castigatus coepi sedulò ipsam inspicere; ac tunc quidem, quae absque indiuisibilibus difficillima visa erat, ipsis opitulantibus, nullo negotio patuit. Modus autem noster ab alijs omnibus quos huc vsque videre contigit, longè diuersus est: et, nisi me nimiùm amo, idem illis omnibus longè antecellit; qui omnium simplicissimus, breuissimus, vniuersalissimus et ad solida detegenda aptissimus existat, vt solus sponte à natura productus, caeteri per vim ab arte effecti videantur. Habes annum quo Trochoidem inuenimus; diem etiam si ità expediret adijcerem: caetera jam ad te scripsi, Et horum omnium testem locupletissimum (praeter quamplurimos alios, quorum epistolas de hac re etiamnum apud me asseruo) ipsum eundem habeo quem laudas, celeberrimum P. Mersennum. Vide ergo num sit cur doleam, cùm vos per exprobrationem obijcitis propositionem illam forsan ante obitum Galilaei nondum fuisse inuentam, qui tamen vixit vsque ad annum 1642. praecipuè, cùm jam ad vos scripserim me anno duodecimo jam elapso inuenisse. Vt sic mihi tot testes habenti, et cui vna sufficere debuit veritas, fidem omnem denegetis. Inuentâ infiniti doctrinâ (liceat adhuc eo nomine vti in hac Epistola; post hac, absit) eáque, pro tempore satis probè excultâ; ego ad tangentes curuarum animum applicui. Ac primùm, vi Analyseos, methodum quandam reperi, quae, etiamsi longè posteà vniuersalis esse deprehensa sit, tamē recèns inuenta, talis non apparuit: quaerebam verò vniuersalem; et particulares methodos (vt adhuc) vbique dedignabar. At Trochoides nostrae occasionem dederunt cur ad motuum compositionem respicerem: occasio satis fuit: ac propositionem vniuersalem tangentium inde deductam vulgauimus circa annum 1636. Extant adhuc, et circumferuntur hac de re lectiones nostrae à nobilissimo D. du Verdus nostro discipulo collectae, atque à multis excriptae. Itaque jamdudùm fide publicâ nobis asserta est talis doctrina, nec alij testes quaerendi, qui omnes habeamus. Circa haec tempora nempe anno 1635, mediante amplissimo senatore Domino de Carcauy, coepi per Epistolas commercium literarum habere cum amplissimo senatore Tholosano Domino de fermat, de quo quid sentiam habes in ea Epistola quam ad R. P. Mersennum direxi super solido vestro hyperbolico infinito. Is ergo vir praestantissimus, primus omnium, duas propositiones nobilissimas ad nos misit sine demonstratione: alteram de parabolis, alteram de planis helicum, vtrisque per omnes dignitatum gradus sumptis (Ne ergo dubites ampliùs, quis primus tales quaestiones proposuerit; illae meae non sunt; quamquam illas ego proprio marte, inuentâ ad id peculiari nostrâ methodo, demonstrauerim) imò vniuersa-
\note{les millésimes 1628, 1634, 1635, 1636, 1642 sont écrits par cette main avec un 1 à boucle d'attaque et un 6 en forme de b (« ıb36 », « ıb42 »). « tamē » : le trait sur l'e (tamen) est celui de la page. « de fermat » est écrit sans majuscule. La dernière syllabe, « -lius », est écrite sous la fin de la ligne, à droite, pour achever « vniuersa- » ; la phrase ne s'arrête pas et se poursuit au-delà de la vue 60.}

\end{document}
